%************************************************
\chapter{Data Management Plan}\label{ch:DataPlan}
%************************************************

This project relies entirely on pre-existing, curated, and de-identified datasets. 
No new primary data collection involving human subjects will be conducted. The data 
management strategy is structured in three progressive phases, transitioning from 
general-purpose computer vision datasets to restricted-access clinical datasets. \\
\\
The research will utilize the following datasets, categorized by project phase: 
\begin{itemize}
    \item \textbf{Phase 1: Prototyping and Baseline Implementation}: Initial architectural 
    development, debugging, and proof-of-concept experiments for disentangled generation 
    and compositionallity analysis will be conducted using well-known, publicly available 
    datasets. These include MNIST, Shapes3D, and CelebA. See Chapter \ref{ch:data} for a 
    detailed list and description of the datasets. These datasets are open-access and 
    require no special data use agreements. 
    \item \textbf{Phase 2: Medical Domain Adaptation}: To transition the generative models to the medical domain, publicly available medical imaging datasets will be utilized. While initial experiments will leverage established repositories like the NIH Chest X-ray and CheXpert datasets, this research is not restricted to radiography. The proposed methods are designed to be adaptable across diverse medical imaging modalities (e.g., MRI, CT, or histopathology). These datasets are fully de-identified and openly available for research purposes, providing a safe and ethical environment to test the models on complex anatomical structures and a wide variety of pathologies. A detailed description of these datasets can be found in Chapter \ref{ch:data}. 
    \item \textbf{Phase 3: Advanced Clinical Application}: For the final evaluation of 
    compositional generation and clinical realism, the project will utilize the MIMIC-CXR dataset. 
    This is a large, restricted-access database of chest radiographs. Access to the MIMIC-CXR 
    dataset has been formally granted via PhysioNet. I have completed the required credentialing 
    (e.g., CITI training on Data or Specimens Only Research) and signed the Data Use Agreement (DUA). 
    No attempts will be made to re-identify any patients. The original restricted data will not 
    be copied to unauthorized personal devices, and access will be limited solely to credentialed 
    researchers involved directly in this project. 
\end{itemize} 
All datasets will be securely stored at the Barcelona Supercomputing Center (BSC - MareNostrum 5) 
and on password-protected, encrypted local workstations used for prototyping. Throughout the project, 
the generative models will produce purely synthetic medical data (e.g., artificial radiographs 
and corresponding segmentation masks or labels). This synthetic data contains no real patient 
information and poses no privacy risks. Derived data, such as model weights, evaluation metrics, 
and latent embeddings, will also be generated and stored securely on the BSC infrastructure. 


%*****************************************
%*****************************************
%*****************************************
%*****************************************
%*****************************************